\chapter*{2024年\textsf{HDU}「高等数学1」期中模拟}

\vspace{-.4\baselineskip}\setcounter{page}{1}
\begin{minipage}{.66\textwidth}
  \centering
  \includegraphics[width = .96\linewidth]{border}
  % \begin{tasks}
  %   [ label-width = 5\ccwd, item-indent = 2cm,
  %     after-item-skip = 0pt ](2)
  %   \task[\textbf{开课学院}：] 理学院
  %   \task[\textbf{解析提供}：] \textsf{秋鱼喵 \& \textsc{A.N. Liu}}
  %   \task[\textbf{课程编码}：] A0714201
  %   \task[\textbf{\textsf\TeX~排版}：] \textsf{Skyrmion}
  % \end{tasks}
\end{minipage}
\hfill
\begin{minipage}{.15\textwidth}
  \centering
  \begin{tikzpicture}
    \node [ inner sep = 0 ] at (0,0) {
      \qrcode[ height = 4em, version = 4 ]{http://weixin.qq.com/r/sRJiemTEGjllrSAZ90ej}};
    \fill [ white ] (-14em/33,14em/33) rectangle (14em/33,-14em/33);
    \fill [ DarkGreen!60, rounded corners = 4em/33 ]
      (-12em/33,12em/33) rectangle (12em/33,-12em/33);
    \node [ inner sep = 0, white, font = \scriptsize ]
      (weixin) at (0,0) {\faIcon{weixin}};
    \node
      [ inner sep = .5ex, align = center, right,
        xshift = 2em, font = \scriptsize ] at (weixin.east) {未\\央\\学\\社};
  \end{tikzpicture}
\end{minipage}
\begin{minipage}{.15\textwidth}
  \centering
  \begin{tikzpicture}
    \node [ inner sep = 0 ] at (0,0) {
      \qrcode [ height = 4em, version = 4 ]{https://qm.qq.com/q/Oo6vNR8ccU}};
    \fill [ white ] (-14em/33,14em/33) rectangle (14em/33,-14em/33);
    \fill [ MidnightBlue!60, rounded corners = 4em/33 ]
      (-12em/33,12em/33) rectangle (12em/33,-12em/33);
    \node [ inner sep = 0, white, font = \scriptsize ]
      (qq) at (0,0) {\faIcon{qq}};
    \node
      [ inner sep = .5ex, align = center, left,
        xshift = -2em, font = \scriptsize ] at (qq.west) {数\\学\\营};
  \end{tikzpicture}
\end{minipage}

\section{选择题（每题3分，共24分）}

\begin{problem}
  设 $y = \sin\frac1x \cos\frac1x \arctan x$，则当 $x \to 0$ 时，$y$ 是
  \begin{tasks}(4)
    \task[\true] 无穷小量
    \task 无穷大量
    \task 有界但非无穷小量
    \task 无界但非无穷大量
  \end{tasks}
\end{problem}
\scratch[4ex]
\begin{solution}
  无穷小乘以有界量为无穷小. $x \to 0$ 时，$\tan^2x \sin\frac1x \sim x^2 \to 0$.
  所以本题选择A项.
\end{solution}

\begin{problem}
  极限 $L = \lim\limits_{x \to 1} \frac{\sqrt{x + 3} - 2}{|x - 1|}$ 的取值为
  \begin{tasks}(4)
    \task $1$
    \task $1/4$
    \task $0$
    \task[\true] 不存在
  \end{tasks}
\end{problem}
\scratch[6ex]
\begin{solution}
  $L = \lim\limits_{x \to 1} \frac{x - 1}{\ab(\sqrt{x + 3} + 2)|x - 1|}$，
  两侧极限分别为 $1/4,~ -1/4$，故极限不存在，所以本题选择D项.
\end{solution}

\begin{problem}
  曲线 $y = x^2$ 与曲线 $y = a\ln x (a \neq 0)$ 相切，则 $a = $
  \begin{tasks}(4)
    \task $4e$
    \task $3e$
    \task[\true] $2e$
    \task $e$
  \end{tasks}
\end{problem}
\scratch[6ex]
\begin{solution}
  对于曲线 $y = x^2 $，其导数为 $\odv y/x = 2x$.
  对于曲线 $y = a \ln x $，其导数为 $\odv y/x = a/x$.
  当 $2x = a/x$，解得相切时 $x = \sqrt{a/2}$，代入两曲线，当 $y$ 值相同时解得 $a = 2e$.
\end{solution}

\begin{problem}
  设 $f'(x_0) = f''(x_0) = 0$，$f'''(x_0) > 0$，则下列选项正确的是
  \begin{tasks}(2)
    \task $f'(x_0)$ 是 $f'(x)$的 极大值
    \task $f(x_0)$ 是 $f(x)$ 的极大值
    \task $f(x_0)$ 是 $f(x)$ 的极小值
    \task[\true] $(x_0,f(x_0))$ 是曲线 $y = f(x)$ 的拐点
  \end{tasks}
\end{problem}
\scratch[14ex]
\begin{solution}\leavevmode
  \begin{itemize}
    \item 由导数定义
    $ f'''(x_0) = \lim\limits_{x \to x_0} \frac{f''(x) - f''(x_0)}{x - x_0}
                = \lim\limits_{x \to x_0} \frac{f''(x)}{x - x_0} > 0$.
    \item 由极限的保号性可知，在去心邻域 $\mathring{U}(x_0,\delta)$ 中
    \begin{itemize}
      \item 当 $x \in (x_0,\delta)$ 时，$x - x_0 > 0,~ f''(x) > 0$.
      \item 当 $x \in (-\delta,x_0)$ 时，$x - x_0 < 0,~ f''(x) < 0$.
    \end{itemize}
    \item 于是，$(x_0,f(x_0))$ 是曲线 $y = f(x)$ 的拐点. 所以本题选择D项.
  \end{itemize}
\end{solution}

\begin{problem}
  设函数 $f(x)$ 在 $x = 0$ 点连续，下列命题中错误的是
  \begin{tasks}(2)
    \task 若 $\lim\limits_{x \to 0} \frac{f(x)}{x}$ 存在，则 $f(0) = 0$
    \task 若 $\lim\limits_{x \to 0} \frac{f(x) + f(-x)}{x}$ 存在，则 $f(0) = 0$
    \task 若 $\lim\limits_{x \to 0} \frac{f(x)}{x}$ 存在，则 $f'(0)$ 存在
    \task[\true] 若 $\lim\limits_{x \to 0} \frac{f(x) + f(-x)}{x}$ 存在，
    则 $f'(0)$ 存在
  \end{tasks}
\end{problem}
\scratch[4ex]
\begin{solution}
  $\sign(x)$ 满足D选项的前提，但是其在 $0$ 点导数不存在.
\end{solution}

\begin{problem}
  设 $f(0) = 0,~ f'(0) = 3$，
  则 $L = \lim\limits_{x \to 0} \frac{f[3\ln(2x + 1)]}{x} = $
  \begin{tasks}(4)
    \task $2$
    \task $6$
    \task $9$
    \task[\true] $18$
  \end{tasks}
\end{problem}
\scratch[10ex]
\begin{solution}
  $ L = \lim\limits_{x \to 0} \frac{f(3\ln(2x + 1)) - f(0)}{x - 0}
      = \lim\limits_{x \to 0}
        \frac{f(3\ln(2x + 1)) - f(0)}{3\ln(2x + 1) - 0} \cdot
        \frac{3\ln(2x + 1)}{x} = f'(3\ln(2x + 1)) \cdot \frac{6x}{x} = 18$.
\end{solution}

\begin{problem}
  设函数 $f(x)$ 连续，且 $f'(0) > 0$，则存在 $\delta > 0$，使得
  \begin{tasks}(2)
    \task $f(x)$ 在 $(0,\delta)$ 内单调增加
    \task $f(x)$ 在 $(-\delta,0)$ 内单调减少
    \task[\true] $\forall x \in (0,\delta)$，有 $f(x) > f(0)$
    \task $\forall x \in (-\delta,0)$，有 $f(x) > f(0)$
  \end{tasks}
\end{problem}
\scratch[4ex]
\begin{solution}
  根据导数的定义，存在区间 $(-\delta,+\delta),~ (f(x) - f(0))/x > 0$.
\end{solution}

\begin{problem}
  设函数 $g(x)$ 可导，$h(x) = e^{1 + g(x)}$，其中 $h'(1) = 1,~ g'(1) = 2$，
  则 $g(1) = $
  \begin{tasks}(4)
    \task $\ln 3 - 1$
    \task $-\ln 3 - 1$
    \task[\true] $-\ln 2 - 1$
    \task $\ln 2 - 1$
  \end{tasks}
\end{problem}
\scratch[6ex]
\begin{solution}
  由题意可知 $h'(x) = g'(x) e^{1 + g(x)}$，
  将 $h'(1) = 1,~ g'(1) = 2$ 代入得 $g(1) = -\ln 2 - 1$.
\end{solution}

\begin{problem}
  使函数 $f(x) = \sqrt[3]{x^2(1 - x)^2}$ 满足罗尔定理条件的区间是
  \begin{tasks}(4)
    \task[\true] $[0,1]$
    \task $[-1,1]$
    \task $[-2,2]$
    \task $[-3/5,4/5]$
  \end{tasks}
\end{problem}
\scratch[4ex]
\begin{solution}
  对于选项BCD中的区间，因函数 $f(x)$ 在 $x = 0$ 处导数不存在，故排除. 所以本题选择A项.
\end{solution}

\begin{problem}
  函数 $f(x) = x^2\sin(x)$，求 $f^{(5)}(0)$ 为
  \begin{tasks}(4)
    \task $20$
    \task[\true] $-20$
    \task $10$
    \task $-10$
  \end{tasks}
\end{problem}
\scratch[8ex]
\begin{solution}
  使用莱布尼兹公式，
  $f^{(5)}(0) = C_5^0 \cdot x^2 \cdot \cos(x) + C_5^1 \cdot 2x \cdot \sin(x)
              + C_5^2 \cdot 2 \cdot (-\cos(x)) + 0 + 0 = -20$.
\end{solution}

\section{填空题（共16分）}

\begin{problem}[本题5分]
  求参数方程的导数：
  $
    \begin{cases}
      x = \cos t\\
      y = \sin t + t
    \end{cases}
  $，求解结果\result{$-\cot t - \csc t$}.
\end{problem}
\scratch[8ex]
\begin{solution}
  $ \odv yx = \frac{\odv y/t}{\odv x/t},~
    \odv xt = -\sin t,~ \odv yt = \cos t + 1$. 于是有
  $y' = \frac{\cos t + 1}{-\sin t} = -\cot t - \csc t$.
\end{solution}

\begin{problem}[本题5分]
  求极限 $L = \lim\limits_{x \to 1} \frac{\cos\frac\pi2 x}{x - 1} = $
  \result{$-\pi/2$}.
\end{problem}
\scratch[8ex]
\begin{solution}
  $ L = \lim\limits_{x \to 1} \frac{\cos\frac\pi2 (x - 1) + \frac\pi2}{x - 1}
      = \lim\limits_{x \to 1} \frac{-\sin\frac\pi2 (x - 1)}{x - 1}
      = \lim\limits_{x \to 1} \frac{-\frac\pi2 (x - 1)}{x - 1}
      = -\frac\pi2$.
\end{solution}

\begin{problem}[本题6分]
  曲线 $y = xe^{-x}$ 的拐点为\result{$(2,2/e^2)$}，
  在拐点处的曲率为 $\kappa = $\result{$0$}.  
\end{problem}
\scratch[12ex]
\begin{solution}
  $y' = e^{-x}(1 - x),~ y'' = e^{-x}(x - 2)$.
  \begin{itemize}
    \item 令 $y'' = 0$，$x = 2$.
      当 $x > 2$ 时，$y'' > 0$；当 $x < 2$ 时，$y'' < 0$. 故 $x = 2$ 为拐点的横坐标.
    \item 当 $x = 2$ 时，$y = 2e^{-2}$. 故拐点为 $(2,2e^{-2})$.
    \item 拐点处的曲率
    $\kappa|_{x = 2} = |y''(2)|/[1 + y'^2(2)]^{3/2}
                     = 0/[1 + (\sim)^2]^{3/2} = 0$.
  \end{itemize}
\end{solution}

\section{解答题（共60分）}

\begin{problem}[本题5分]
  求极限 $L = \lim\limits_{{x \to \infty}} (1^x + 2^x + 3^x)^{1/x}$.
\end{problem}
\scratch[12ex]
\begin{solution}
  $ L = \lim\limits_{{x \to \infty}} (1^x + 2^x + 3^x)^{1/x}
      = \lim\limits_{{x \to \infty}}
        (3^x)^{1/x}\ab[(\frac13)^x + (\frac23)^x + 1^x]^{1/x}
      = \lim\limits_{{x \to \infty}} 3(0 + 0 + 1)^{1/x} = 3$.
\end{solution}

\begin{problem}[本题5分]
  求极限：
  $L = \lim\limits_{x \to 2} \frac{\sqrt{5x - 1} - \sqrt{2x + 5}}{x^2 - 4}$.
\end{problem}
\scratch[18ex]
\begin{solution}
  将分子有理化
  \[
    L = \lim_{x \to 2}
        \frac{\ab(\sqrt{5x - 1} - \sqrt{2x + 5})
              \ab(\sqrt{5x - 1} + \sqrt{2x + 5})}
          {(x^2 - 4)\ab(\sqrt{5x - 1} + \sqrt{2x + 5})}
      = \lim_{x \to 2}
        \frac{(5x - 1) - (2x + 5)}{(x^2 - 4) \ab(\sqrt9 + \sqrt9)}
      = \lim_{x \to 2} \frac{3(x - 2)}{6(x - 2)(x + 2)} = \frac18
  \]
\end{solution}

\begin{problem}[本题5分]
  函数 $y = y(x)$ 由方程 $\ln(x^2 + y) = x^3y + \sin x$ 确定，
  则 $\odv yx \bigg|_{x = 0}$ 为？
\end{problem}
\scratch[20ex]
\begin{solution}\leavevmode
  \begin{itemize}
    \item 将 $x = 0$ 代入方程：
    $\ln(0^2 + y) = 0^3 \cdot y + \sin 0,~ \ln y = 0,~ y = 1$.
    \item 对方程两边同时求导：
    $\frac{1}{x^2 + y} \cdot (2x + y') = 3x^2y + x^3y' + \cos x$.
    \item 将 $x = 0,~ y = 1$ 代入上述方程
    $ \frac{1}{0^2 + 1} \cdot (2 \cdot 0 + y'(0))
    = 3 \cdot 0^2 \cdot 1 + 0^3 \cdot y' + \cos 0,~
      \odv yx\Big|_{x = 0} = y'(0) =  1$.
  \end{itemize}
\end{solution}

\begin{problem}[本题5分]
  $y = x^{x^x}$，求$y'$.
\end{problem}
\scratch[16ex]
\begin{solution}\leavevmode
  \begin{itemize}
    \item 对 $y = x^{x^x}$ 两侧取对数
    $
      \ln y = \ln x^{x^x} = x^x\ln x,~
      \ln\ln y = \ln(x^x\ln x ) = \ln(x^x) + \ln\ln x = x\ln x + \ln\ln x
    $.
    \item 两侧求导
    $\frac{1}{\ln y} \cdot \frac1y \cdot y' = \ln x + 1 + \frac{1}{x\ln x}$，
    整理得 $y' = x^{x^x + x - 1} (x\ln ^2x + x\ln x + 1)$.
  \end{itemize}
\end{solution}

\begin{problem}[本题6分]
  设 $y = \frac{e^x + 1}{e^x - 1}$，求
  \begin{tasks}[ label = \arabic*. ](2)
    \task 函数的间断点并判断其类型.
    \task 该函数的水平渐近线和垂直渐近线.
  \end{tasks}
\end{problem}
\scratch[24ex]
\begin{solution}\leavevmode
  \begin{enumerate}
    \item 当 $x = 0$ 时，函数无定义，因此 $x = 0$ 为间断点.
    \item 当 $x = 0$ 时，函数无定义，
    $\lim\limits_{x \to 0} \frac{e^x + 1}{e^x - 1} = \infty$，
    因此 $x = 0$ 为垂直渐近线.
    \[
        \lim_{x \to +\infty} f(x) = \lim_{x \to +\infty} \frac{e^x + 1}{e^x - 1}
      = \lim_{x \to +\infty} \frac{1 + \frac{1}{e^x}}{1 - \frac{1}{e^x}}
      = \frac{1 + 0}{1 - 0} = 1,~
        \lim_{x \to -\infty} f(x) = \lim_{x \to -\infty} \frac{e^x + 1}{e^x - 1}
      = \frac{0 + 1}{0 - 1} = -1
    \]
    因此，$y = -1,~ y = 1$为水平渐近线.
  \end{enumerate}
\end{solution}

\begin{problem}[本题7分]
  多项式 $P(x) = a_1x + a_2x^2 + \cdots + a_nx^n$ 的系数满足
  $\sum^n_{i = 1}a_i = 0$.
  证明在 $(0,2)$ 上多项式 $a_1t + 2a_2t^2 + \cdots + na_nt^n = 0$ 有实数解.
\end{problem}
\scratch[10ex]
\begin{solution}
  \begin{proof}
    显然 $P(0) = P(1) = 0$，将第二个多项式除以 $t$ 得
    $a_1 + 2a_2t + \cdots + na_nt^{n - 1} = 0$，
    后者刚好为 $P'(t)$. 根据罗尔定理即得证.
  \end{proof}
\end{solution}

\begin{problem}[本题5分]
  求极限 $L = \lim\limits_{{x \to \infty}}  x - x^2\ln\ab(1 + \frac1x)$.
\end{problem}
\scratch[14ex]
\begin{solution}
  令 $t = 1/x$，使用洛必达法则
  \[
    L = \lim_{t \to 0} \frac1t - \frac{1}{t^2}\ln(1 + t)
      = \lim_{t \to 0} \frac{t - \ln(1 + t)}{t^2}
      \xlongequal{\text{L'Hôpital}} \lim_{t \to 0} \frac{1 - 1/(1 + t)}{2t}
      = \lim_{t \to 0} \frac{1/(1 + t)^2}{2} = 1/2
  \]
\end{solution}

\begin{problem}[本题5分]
  设 $y = f(\ln x)e^{f(x)}$，其中 $f$ 可微，则 $\d y$ 为？
\end{problem}
\scratch[8ex]
\begin{solution}
  $
     \d y = y'\d x = [f(\ln x)]' \cdot e^{f(x)} + f(\ln x)\ab[e^{f(x)}]'\d x
          = e^{f(x)} \ab[\frac{f'(\ln x)}{x} + f(\ln x)f'(x)] \d x  
  $.
\end{solution}

\begin{problem}[本题6分]
  讨论函数 $f(x) =
  \begin{cases}
    \frac{x}{1 + e^{\frac1x}}, & x\neq 0\\
    0,                         & x = 0
  \end{cases}$ 在点 $x = 0$ 的连续性和可导性质.
\end{problem}
\scratch[28ex]
\begin{solution}\leavevmode
  \begin{itemize}
    \item 连续性讨论左右极限和函数值
    \[
      f(0^-) = \lim_{x \to 0^-}\frac{x}{1 + e^{1/x}} = \frac{0}{1 + 0} = 0,~
      f(0) = 0,~
      f(0^-) = \lim_{x \to 0^+}\frac{x}{1 + e^{1/x}} = \frac{0}{1 + \infty} = 0
    \]
    于是，$f(0^-) = f(0) = f(0^+)$，$f(x)$ 在 $x = 0$ 处连续.
    \item 可导性讨论左右导数
    \begin{gather*}
      f'(0^-) = \lim_{x \to 0^-} \frac{f(x) - f(0)}{x - 0}
              = \lim_{x \to 0^-} \frac{\frac{x}{1 + e^{1/x}} - 0}{x}
              = \lim_{x \to 0^-} \frac{x}{(1 + e^{1/x})}
              = \lim_{x \to 0^-} \frac{1}{1 + e^{1/x}} = \frac{1}{1 + 0} = 1\\
      f'(0^+) = \lim_{x \to 0^+} \frac{f(x) - f(0)}{x-0}
              = \lim_{x \to 0^+} \frac{\frac{x}{1 + e^{1/x}} - 0}{x}
              = \lim_{x \to 0^+} \frac{x}{(1 + e^{1/x})}
              = \lim_{x \to 0^+} \frac{1}{1 + e^{1/x}}
              = \frac{1}{1 + \infty} = 0
    \end{gather*}
    于是，$f'(0^-) \neq f'(0^+)$，$f(x)$ 在 $x = 0$ 处不可导.
  \end{itemize}
\end{solution}